% =====================================================================
%  3.3 Polynomfunksjoner — fasit
% =====================================================================

\fk{Del 1 — Polynomdivisjon og faktorisering}
\fo{1.1} a) $x + 5$ \ b) $x + 3$ \ c) $2x + 3$ (alle uten rest)
\fo{1.2} a) Kvotient $x^2 + 4x + 3$, rest $0$. \ b) $(x-2)(x^2 + 4x + 3) = x^3 + 2x^2 - 5x - 6$
\fo{1.3} a) $P(3) = 0$: ja \ b) $P(-2) = -20$: nei \ c) $P(-1) = 0$: ja
\fo{1.4} a) $x^2 + 3x - 4$ \ b) $x^2 - 2x + 4$ \ c) $2x^2 - 5x - 3$ \ d) $x^2 + 2x + 3$ rest $7$, altså $x^2 + 2x + 3 + \frac{7}{x-1}$
\fo{1.5} a) $(x-1)(x-2)(x-3)$ \ b) $(x-2)(x+2)(x+3)$ \ c) $(x+1)^2(x-2)$ (divisjon med $x + 1$ og en tom kolonne for $x^2$)
\fo{1.6} Feilen er i andre runde: $-3x(x-2) = -3x^2 + 6x$, og $-3x^2 + 2x$ pluss det motsatte, $3x^2 - 6x$, gir $-4x$, ikke $4x$. Faktoren $(x-2)$ var riktig. Riktig kvotient er $x^2 - 3x - 4$, så $P(x) = (x-2)(x-4)(x+1)$.
\fo{1.7} a) Kvotienten er det hele tallet $13$ i $157 : 12$, og her polynomet $x^2 + 2x + 3$. Resten er $1$ i tallet og her $7$. \ b) Når resten har lavere grad enn divisoren, kan ikke ledende ledd i resten deles på ledende ledd i divisoren. Da er det ikke mer å dele, på samme måte som $1 < 12$ i $157 : 12$.
\fo{1.8} a) Noen ganger. $x^3 - x$ har tre nullpunkter ($0$, $1$ og $-1$), men $x^3 + 2x$ har bare $x = 0$.
b) Noen ganger. Resten forskjellig fra $0$ viser bare at $a$ ikke er et nullpunkt. $x^3 - 4x$ gir rest $-3$ ved $(x-1)$ og har likevel tre nullpunkter, mens $x^2 + 1$ gir rest $2$ og har ingen.
c) Alltid. $P(2) = 0$ gir $P(x) = (x-2)Q(x)$. Så er $P(3) = 1 \cdot Q(3) = 0$, og da er $(x-3)$ en faktor i $Q$.
\fo{1.9} a) $P(2) = 8 + 4k - 8 + 4 = 4 + 4k = 0$, så $k = -1$. \ b) $x^3 - x^2 - 4x + 4 = (x-2)(x^2 + x - 2) = (x-1)(x-2)(x+2)$
\fo{1.10} Flere svar er riktige. Eksempel: $P(x) = x^3 - 3$. Resten er $P(2) = 8 - 3 = 5$ fordi $P(x) = (x-2)Q(x) + r$ gir $P(2) = r$. Ja, mange. Et annet eksempel er $x^3 + x - 5$, for $P(2) = 8 + 2 - 5 = 5$.
\fo{Sjekk 1} a) $x + 5$ \ b) Ja, $P(2) = 8 - 12 + 4 = 0$. \ c) Divisjon med rest betyr nettopp $P(x) = D \cdot Q + r$. Når vi ganger ut høyresiden, må vi få $P(x)$ tilbake. Gjør vi ikke det, er det en regnefeil.

\fk{Del 2 — Polynomfunksjoner: nullpunkter og forløp}
\fo{2.1} a) Grad $3$, ledende ledd $2x^3$. Opp til høyre, ned til venstre. \ b) Grad $4$, $-x^4$. Ned på begge sider. \ c) Grad $3$, $-x^3$. Ned til høyre, opp til venstre.
\fo{2.2} a) $x = -4$, $1$, $5$ og $(0, 20)$ \ b) $x = -2$, $0$, $3$ og $(0, 0)$ \ c) $x = -5$, $2$, $6$ og $(0, -60)$
\fo{2.3} a) Produktlinja er stiplet til venstre for $-1$, heltrukket mellom $-1$ og $2$, stiplet mellom $2$ og $4$ og heltrukket til høyre for $4$. \ b) $x \in \langle \leftarrow, -1\rangle$ eller $x \in \langle 2, 4\rangle$
\fo{2.4} a) $(x-2)(x-4)(x+3) = 0$: $x = -3$, $2$ eller $4$ \ b) $x(x+2)^2 = 0$: $x = 0$ eller $x = -2$ \ c) $x(x-4)(x+4) = 0$: $x = -4$, $0$ eller $4$
\fo{2.5} a) $(x-2)(x-5)(x+3) < 0$: $\langle \leftarrow, -3\rangle$ og $\langle 2, 5\rangle$ \ b) $(x-2)^2(x+4) > 0$: $\langle -4, 2\rangle$ og $\langle 2, \to\rangle$ \ c) $(x+1)(x-3)^2 \ge 0$. Faktoren $(x-3)^2$ er $\ge 0$ og er $0$ i $x = 3$. Ellers bestemmer $(x+1)$ fortegnet, så løsningen er $x \ge -1$, altså $[-1, \to\rangle$.
\fo{2.6} a) Nullpunktene er $x = -3$ og $x = 1$ (dobbelt). $f(0) = -3$. Ledende ledd er $-x^3$: opp til venstre, ned til høyre. \ b) Grafen kommer ovenfra, skjærer $x$-aksen i $-3$, går under aksen, berører den i $x = 1$ og går videre nedover.
\fo{2.7} Feilen er steget «x = 1 tas med». Ulikheten er streng, og produktet er $0$ i $x = 1$, ikke mindre enn $0$. Riktig svar er $x \in \langle \leftarrow, -3\rangle$ (også $x = -3$ er utenfor). Kontroll: $x = -4$ gir $-25$ og $x = 1$ gir $0$.
\fo{2.8} a) Alltid. Den ene enden er under $x$-aksen og den andre over (for både $a_n > 0$ og $a_n < 0$). Grafen har ingen brudd, så den må krysse. \ b) Noen ganger. $x^4 + 1$ har ingen nullpunkter, $x^4 - 1$ har to. \ c) Noen ganger. $x^2(x^2+1) \ge 0$, men $(x-1)^2(x-2)$ er negativ for $x = \frac{3}{2}$.
\fo{2.9} a) $f(x) = a(x+1)(x-3)(x-4)$ og $f(0) = 12a = -24$, så $a = -2$: $f(x) = -2(x+1)(x-3)(x-4)$. Det finnes bare én: tre nullpunkter fastlegger faktorene, og $f(0)$ bestemmer tallet foran. \ b) Ja, med ett dobbelt nullpunkt. $-2(x+1)^2(x-3)(x-4)$, $\frac{2}{3}(x+1)(x-3)^2(x-4)$ og $\frac{1}{2}(x+1)(x-3)(x-4)^2$ har alle $f(0) = -24$.
\fo{2.10} a) $f(0) = 0$, så $x = 0$ er et nullpunkt. Da må $f(x) = x(x-2)^2 = x^3 - 4x^2 + 4x$, så $a = -4$ og $b = 4$. \ b) $f(x) = x(x-2)^2 \le 0$ gir $x \le 0$ eller $x = 2$, altså $\langle \leftarrow, 0]$ og $x = 2$. For $x \neq 2$ er $(x-2)^2 > 0$, så fortegnet bestemmes av $x$. I $x = 2$ er $f(2) = 0$, og $\le$ tillater $0$.
\fo{Sjekk 2} a) Grad $4$, ledende ledd $-2x^4$. \ b) $x = 0$ og $x = 3$ \ c) For $x = 100$ er $x^4 = 10^8$, mens leddene av lavere grad er mye mindre. Ledende ledd dominerer.

\fk{Ekstra utfordring}
\fo{U1} $(x^3 + 0x^2 + 0x + 2) : (x - 1) = x^2 + x + 1$ rest $3$. Altså $\dfrac{x^3+2}{x-1} = x^2 + x + 1 + \dfrac{3}{x-1}$. Kontroll $x = 3$: $\frac{29}{2} = 14,5$ og $9 + 3 + 1 + \frac{3}{2} = 14,5$.
\fo{U2} a) $x > 0$, $20 - 2x > 0$ og $30 - 2x > 0$ gir $0 < x < 10$. \ b) $V(x) - 1000 = 4(x-5)(x^2 - 20x + 50)$. Løsninger: $x = 5$ og $x \approx 2,93$. Den tredje, $x \approx 17,07$, er utenfor definisjonsmengden.
\fo{U3} a) $-2$, $-1$ og $4$, hver på egen linje. $P(x) = (x+1)(x+2)(x-4)$. \ b) Bytt ut funksjonen med \lstinline[style=laereverk,numbers=none,frame=none]|return x**4 - 10*x**2 + 9| og sett $d = 9$. Utskrift: $-3$, $-1$, $1$, $3$.
\fo{U4} $x^4 - 5x^2 + 4 = (x-1)(x^3 + x^2 - 4x - 4) = (x-1)(x+1)(x-2)(x+2)$. Løsninger: $x = -2$, $-1$, $1$, $2$.
\fo{U5} a) Fra $a^3 + ba^2 + ca + d = 0$ får vi $d = -a(a^2 + ba + c)$. Parentesen er et heltall, så $a$ går opp i $d$. \ b) Et heltallig nullpunkt må gå opp i $d$, så ingen andre heltall kan være nullpunkter. \ c) Nei. For $x^3 - 2$ går $2$ opp i $-2$, men $P(2) = 6 \neq 0$.
\fo{U6} a) $P(1) = 1 + a + b = 5$ og $P(-1) = -1 - a + b = -1$. Legger vi sammen, får vi $2b = 4$, så $b = 2$ og $a = 2$. \ b) Resten er $P(2) = 8 + 4 + 2 = 14$.
\fo{U7} a) Partall $n$: grafen berører $x$-aksen og snur, ingen fortegnsskifte. Oddetall $n$: grafen skjærer aksen og skifter fortegn. \ b) $(x-1)^n$ har fortegnet til $(x-1)$ opphøyd i $n$. Et partall faktorer gir alltid positivt produkt, også når $x < 1$. Et oddetall beholder fortegnet. \ c) Flere svar er riktige. Eksempel: $f(x) = (x+2)(x-1)^4$ (grad 5, ledende ledd $x^5$, ned til venstre, opp til høyre).
